% ejemplo_uso.tex — Ejemplos del sistema visual "Vía Accesible" con la plantilla TFUOC
% Copia en la carpeta de la plantilla (junto a TFUOC.cls): tfm-estilo.sty
% Compila: pdflatex → biber → pdflatex → pdflatex (Overleaf lo hace solo).
\documentclass[IB]{TFUOC}
\usepackage[tabla,bibliografia]{tfm-estilo}   % quita [tabla] si prefieres «Cuadro»; añade final para la entrega
\addbibresource{referencias.bib}

\title{Clasificación de la accesibilidad de calles mediante segmentación semántica de imágenes aéreas}
\titcrt{Accesibilidad de calles desde el aire}
\author{Víctor Pariente}
\date{2026}
\nomPDC{Nombre Tutor}
\nomPRA{Nombre Profesorado Responsable}
\titulac{Máster en Ciencia de Datos}
\area{Área del trabajo}
\idioma{Castellano}
\credits{15}
\parcla{segmentación semántica, accesibilidad, imágenes aéreas}
\licenc{ccBy}
\abstractidioma{Resumen.}
\abstractenglish{Abstract.}

\begin{document}
\estructura
\tableofcontents
\listoffigures
\listoftables

\chapter{Ejemplos del sistema visual}

\section{Gráfico de líneas}
La figura~\ref{fig:perdida} muestra la pérdida por época.

\begin{figure}[htbp]
  \centering
  \begin{tikzpicture}
    \begin{axis}[tfm, width=0.8\textwidth, height=0.45\textwidth,
        xlabel={Época}, ylabel={Pérdida (u.a.)}, legend pos=north east]
      \addplot coordinates {(1,0.92)(2,0.66)(3,0.52)(4,0.44)(5,0.39)(6,0.35)(7,0.33)(8,0.31)};
      \addlegendentry{Entrenamiento}
      \addplot coordinates {(1,0.98)(2,0.74)(3,0.61)(4,0.55)(5,0.51)(6,0.49)(7,0.48)(8,0.48)};
      \addlegendentry{Validación}
      \addplot[serie-base, dashed, line width=\trazo] coordinates {(1,0.45)(8,0.45)};
      \addlegendentry{Referencia (FCN)}
    \end{axis}
  \end{tikzpicture}
  \caption{Pérdida de entrenamiento y validación por época (datos de ejemplo).}
  \label{fig:perdida}
  \fuente{elaboración propia.}
\end{figure}

\section{Barras agrupadas}
\begin{figure}[htbp]
  \centering
  \begin{tikzpicture}
    \begin{axis}[tfmbarras, width=0.9\textwidth, height=0.45\textwidth,
        ylabel={IoU}, xlabel={Clase}, ymin=0, ymax=1,
        symbolic x coords={Calzada,Acera,Paso peat.,Vegetación,Obstáculo},
        legend style={at={(0.5,1.03)}, anchor=south}, legend columns=3]
      \addplot coordinates {(Calzada,0.91)(Acera,0.68)(Paso peat.,0.58)(Vegetación,0.84)(Obstáculo,0.41)};
      \addplot coordinates {(Calzada,0.92)(Acera,0.69)(Paso peat.,0.62)(Vegetación,0.85)(Obstáculo,0.44)};
      \addplot coordinates {(Calzada,0.93)(Acera,0.73)(Paso peat.,0.66)(Vegetación,0.86)(Obstáculo,0.49)};
      \legend{U-Net, DeepLabV3+, SegFormer-B2}
    \end{axis}
  \end{tikzpicture}
  \caption{IoU por clase para cada modelo (datos de ejemplo).}
\end{figure}

\section{Tabla de resultados}
\begin{table}[htbp]
  \centering
  \caption{Resultados de segmentación en el conjunto de validación (datos de ejemplo).}
  \begin{tabular}{lrrr}
    \toprule
    \cabeceratabla Modelo & mIoU & IoU acera & F1 \\
    \midrule
    FCN-8s (referencia) & 0,61 & 0,55 & 0,70 \\
    U-Net (ResNet-34)   & 0,72 & 0,68 & 0,79 \\
    SegFormer-B2        & \mejor{0,77} & \mejor{0,73} & \mejor{0,83} \\
    \bottomrule
  \end{tabular}
\end{table}

\section{Escala de accesibilidad}
\claseacc{accesible}{Accesible}\quad
\claseacc{practicable}{Practicable}\quad
\claseacc{deficiente}{Deficiente}\quad
\claseacc{no-accesible}{No accesible}\quad
\claseacc{sin-datos}{Sin datos}

\section{Diagrama de red (U-Net simplificada)}
\begin{figure}[htbp]
  \centering
  \begin{tikzpicture}[tfm, node distance=3mm]
    \node[nnentrada, minimum height=22mm] (in) {RGB\\$512^2$};
    \node[nnconv, right=of in, minimum height=22mm] (e1) {64};
    \node[nnpool, right=of e1, minimum height=14mm, minimum width=4mm] (p1) {};
    \node[nnconv, right=of p1, minimum height=14mm] (e2) {128};
    \node[nnatencion, right=of e2, minimum height=9mm] (b) {256};
    \node[nnup, right=of b, minimum height=14mm, minimum width=4mm] (u1) {};
    \node[nnconv, right=of u1, minimum height=14mm] (d2) {128};
    \node[nnup, right=of d2, minimum height=22mm, minimum width=4mm] (u2) {};
    \node[nnconv, right=of u2, minimum height=22mm] (d1) {64};
    \node[nnsalida, right=of d1, minimum height=22mm] (out) {Softmax};
    \foreach \a/\b in {in/e1, e1/p1, p1/e2, e2/b, b/u1, u1/d2, d2/u2, u2/d1, d1/out}
      \draw[flecha] (\a) -- (\b);
    \draw[skip] (e1.north) to[bend left=25] (d1.north);
    \draw[skip] (e2.north) to[bend left=30] (d2.north);
  \end{tikzpicture}
  \caption{Arquitectura encoder–decoder con conexiones \emph{skip}.}
\end{figure}

\section{Diagrama de arquitectura}
\begin{figure}[htbp]
  \centering
  \begin{tikzpicture}[tfm, node distance=5mm]
    \node[arqexterno] (pnoa) {Ortofoto\\PNOA};
    \node[arqmodulo, right=of pnoa] (tes) {Teselado};
    \node[arqdatos, right=of tes] (tiles) {Teselas};
    \node[arqmodulo, right=of tiles] (seg) {Modelo de\\segmentación};
    \node[arqdatos, right=of seg] (mask) {Máscaras};
    \node[arqmodulo, below=8mm of mask] (cla) {Clasificador de\\accesibilidad};
    \node[arqexterno, left=of cla] (sig) {SIG\\municipal};
    \draw[flecha] (pnoa) -- (tes); \draw[flecha] (tes) -- (tiles);
    \draw[flecha] (tiles) -- (seg); \draw[flecha] (seg) -- (mask);
    \draw[flecha] (mask) -- (cla); \draw[flecha] (cla) -- (sig);
    \begin{scope}[on background layer]
      \node[arqgrupo, fit=(tes)(mask)(cla), label={[etiqueta]above left:Pipeline en Python}] {};
    \end{scope}
  \end{tikzpicture}
  \caption{Arquitectura del sistema.}
\end{figure}

\section{Imagen externa}
La figura~\ref{fig:externa} usa una imagen de la plantilla como ejemplo; sustitúyela por tu ortofoto.

\figuraexterna[0.45\textwidth]{Rplotmanh.png}{Ejemplo de imagen externa con su fuente.}{fig:externa}{plantilla TFM de la UOC.}

\begin{figure}[htbp]
  \centering
  \begin{subfigure}{0.3\textwidth}\includegraphics[width=\linewidth]{Rplotmanh.png}\caption{Imagen}\end{subfigure}\hfill
  \begin{subfigure}{0.3\textwidth}\includegraphics[width=\linewidth]{Rplotmanh.png}\caption{Referencia}\end{subfigure}\hfill
  \begin{subfigure}{0.3\textwidth}\includegraphics[width=\linewidth]{Rplotmanh.png}\caption{Predicción}\end{subfigure}
  \caption{Mosaico imagen, referencia y predicción.}
\end{figure}

\section{Texto a dos columnas}
El texto normal de la memoria sigue a una columna. Los bloques que lo necesiten (resumen ampliado, estado del arte, comparativas) pueden ir a dos o más columnas.

\begin{doscolumnas}[filete=true]
Los planes de accesibilidad municipales requieren inventariar cada tramo de calle: ancho libre de paso, pendientes, rebajes en los pasos de peatones y obstáculos. La inspección in situ es lenta y costosa, y rara vez se actualiza.

Las ortofotos de alta resolución permiten observar gran parte de estos elementos desde el aire. Un modelo de segmentación semántica separa calzada, acera, pasos de peatones, vegetación y obstáculos, y a partir de las máscaras se calculan métricas por tramo.

\begin{figuracolumna}
  \begin{tikzpicture}
    \begin{axis}[tfm, height=4.5cm, xlabel={Época}, ylabel={mIoU}]
      \addplot coordinates {(1,0.41)(2,0.55)(3,0.63)(4,0.68)(5,0.71)};
    \end{axis}
  \end{tikzpicture}
  \caption{mIoU en validación (datos de ejemplo).}
  \label{fig:columna}
\end{figuracolumna}

La figura~\ref{fig:columna} se compone dentro de la columna, sin flotar, y se numera con las demás figuras. Las tablas usan \texttt{tablacolumna} de la misma forma.

\begin{tablacolumna}
  \caption{Resultados (ejemplo).}
  \begin{tabular}{lr}
    \toprule
    \cabeceratabla Modelo & mIoU \\
    \midrule
    U-Net & 0,72 \\
    SegFormer-B2 & \mejor{0,77} \\
    \bottomrule
  \end{tabular}
\end{tablacolumna}

Con \texttt{equilibrar=false} la primera columna se llena entera antes de pasar a la siguiente; con \texttt{columnas=3} el bloque se reparte en tres.
\end{doscolumnas}

\begin{doscolumnas}[columnas=3, separacion=6mm, titulo={\subsection*{Tres columnas con título común}}]
\textbf{Calzada.} Superficie rodada; delimita el ancho disponible para la acera.

\textbf{Acera.} Itinerario peatonal; su ancho libre define la clase de accesibilidad.

\textbf{Obstáculos.} Mobiliario y vehículos que reducen el paso efectivo.
\end{doscolumnas}

\section{Diagrama de flujo}
\begin{figure}[htbp]
  \centering
  \begin{tikzpicture}[tfm, node distance=5mm and 7mm]
    \node[flujoterminal] (ini) {Inicio};
    \node[flujodatos, below=of ini] (orto) {Ortofoto PNOA};
    \node[flujoalgoritmo, below=of orto] (seg) {Segmentación\\semántica (SegFormer)};
    \node[flujodecision, below=of seg] (conf) {¿Confianza\\$\geq 0{,}8$?};
    \node[flujoalgoritmo, below=of conf] (vec) {Vectorización\\y métricas por tramo};
    \node[flujoproceso, right=12mm of conf] (rev) {Revisión\\manual};
    \node[flujoalgoritmo, below=of vec] (cla) {Clasificación de\\accesibilidad};
    \node[flujodatos, below=of cla] (sal) {Capa SIG por tramo};
    \node[flujoterminal, below=of sal] (fin) {Fin};
    \draw[flecha] (ini) -- (orto); \draw[flecha] (orto) -- (seg); \draw[flecha] (seg) -- (conf);
    \draw[flecha] (conf) -- node[flujoetiqueta, right] {sí} (vec);
    \draw[flecha] (conf) -- node[flujoetiqueta, above] {no} (rev);
    \draw[flecha] (rev) |- (vec);
    \draw[flecha] (vec) -- (cla); \draw[flecha] (cla) -- (sal); \draw[flecha] (sal) -- (fin);
  \end{tikzpicture}
  \caption{Flujo del procesado: la salida de cada algoritmo alimenta al siguiente.}
\end{figure}

\section{Esquema conceptual}
\begin{figure}[htbp]
  \centering
  \begin{tikzpicture}[tfm]
    \node[conceptocentral] (c) {Accesibilidad\\de la calle};
    \node[concepto] (a) at (-4.2,1.4) {Ancho libre\\de acera};
    \node[concepto] (p) at (4.2,1.4) {Pasos de\\peatones};
    \node[concepto] (o) at (-4.2,-1.4) {Obstáculos};
    \node[concepto] (e) at (4.2,-1.4) {Pendiente};
    \node[conceptosec] (m) at (0,-2.6) {Máscara de segmentación};
    \foreach \n/\t in {a/determina, p/condiciona, o/reduce, e/limita}
      \draw[relacion] (\n) -- node[relacionetiqueta] {\t} (c);
    \draw[relacion] (m) -- node[relacionetiqueta] {mide} (o);
    \draw[relacion] (m) -- node[relacionetiqueta, pos=0.4] {no mide} (e);
  \end{tikzpicture}
  \caption{Esquema conceptual de los factores de accesibilidad.}
\end{figure}

\section{Estados}
Estados en línea: \estado{pendiente} \estado{en-curso} \estado{revision} \estado{completado} \estado{bloqueado} \estado[Validado]{completado}.

\estadobloque{revision}{Revisar las cifras de la tabla 3.2 con el tutor.}

\begin{table}[htbp]
  \centering
  \caption{Estado de los objetivos.}
  \begin{tabular}{llc}
    \toprule
    \cabeceratabla Objetivo & Estado & \\
    \midrule
    O1. Dataset anotado & \estado{completado} & \glifoestado{completado} \\
    O2. Modelo de segmentación & \estado{revision} & \glifoestado{revision} \\
    O3. Clasificación por tramo & \estado{en-curso} & \glifoestado{en-curso} \\
    O4. Validación en campo & \estado{bloqueado} & \glifoestado{bloqueado} \\
    \bottomrule
  \end{tabular}
\end{table}

\section{Tablero Kanban}
\begin{kanban}[4]
  \kanbancolumna{Pendiente}{\tarjeta{Validación en campo}{pendiente}{PEC 4}\tarjeta{Anexo de instalación}{pendiente}{}}
  \kanbancolumna{En curso}{\tarjeta{Clasificador por tramo}{en-curso}{Reglas de ancho libre}}
  \kanbancolumna{En revisión}{\tarjeta{Entrenamiento SegFormer}{revision}{Tutor}}
  \kanbancolumna{Hecho}{\tarjeta{Dataset anotado}{completado}{1\,200 teselas}\tarjeta{Descarga PNOA}{completado}{}}
\end{kanban}

\section{Recuadros}
\begin{recuadro}[definicion]{Itinerario peatonal accesible}
Ámbito de paso peatonal con ancho libre mínimo de 1,80~m y sin escalones.
\end{recuadro}
\begin{recuadro}[atencion]{}
Las sombras de edificios altos reducen la calidad de la segmentación.
\end{recuadro}
\begin{recuadro}[hipotesis]{H1}
La segmentación permite estimar el ancho libre con un error inferior a 0,3~m.
\end{recuadro}

\section{Diagrama de secuencia}
\begin{figure}[htbp]
  \centering
  \begin{tikzpicture}[tfm]
    \foreach \x/\n/\t/\s in {0/u/Técnico/participante, 3.6/p/Pipeline/participante, 7.2/m/Modelo/participante, 10.8/g/SIG municipal/participanteext}
      {\node[\s] (\n) at (\x,0) {\t}; \draw[lineavida] (\n.south) -- ++(0,-5.2);}
    \fill[activacion] (3.5,-0.9) rectangle (3.7,-4.6);
    \fill[activacion] (7.1,-1.6) rectangle (7.3,-2.6);
    \draw[activacion] (3.5,-0.9) rectangle (3.7,-4.6); \draw[activacion] (7.1,-1.6) rectangle (7.3,-2.6);
    \draw[mensaje] (0,-1) -- node[mensajetexto] {1. procesar(zona)} (3.5,-1);
    \draw[mensaje] (3.7,-1.7) -- node[mensajetexto] {2. predecir(teselas)} (7.1,-1.7);
    \draw[respuesta] (7.1,-2.5) -- node[mensajetexto] {3. máscaras} (3.7,-2.5);
    \draw[mensaje] (3.7,-3.4) -- node[mensajetexto] {4. publicar(capa)} (10.8,-3.4);
    \draw[respuesta] (3.5,-4.4) -- node[mensajetexto] {5. informe} (0,-4.4);
  \end{tikzpicture}
  \caption{Secuencia de una ejecución del sistema.}
\end{figure}

\section{Barras horizontales, apiladas y líneas verticales}
\begin{figure}[htbp]
  \centering
  \begin{subfigure}{0.48\textwidth}
  \begin{tikzpicture}
    \begin{axis}[tfmbarrash, width=\linewidth, height=6cm, xlabel={IoU}, xmin=0, xmax=1,
        symbolic y coords={Calzada,Acera,Paso peat.,Obstáculo}, legend pos=south east]
      \addplot coordinates {(0.91,Calzada)(0.68,Acera)(0.58,Paso peat.)(0.41,Obstáculo)};
      \addplot coordinates {(0.93,Calzada)(0.73,Acera)(0.66,Paso peat.)(0.49,Obstáculo)};
      \legend{U-Net, SegFormer-B2}
    \end{axis}
  \end{tikzpicture}\caption{Horizontal agrupado}
  \end{subfigure}\hfill
  \begin{subfigure}{0.48\textwidth}
  \begin{tikzpicture}
    \begin{axis}[tfmapiladash, width=\linewidth, height=6cm, xlabel={Tramos (\%)}, xmin=0, xmax=100,
        cycle list name=tfmaccesibilidad, symbolic y coords={Centro,Ensanche,Periferia}]
      \addplot coordinates {(52,Centro)(38,Ensanche)(25,Periferia)};
      \addplot coordinates {(24,Centro)(30,Ensanche)(28,Periferia)};
      \addplot coordinates {(15,Centro)(20,Ensanche)(27,Periferia)};
      \addplot coordinates {(7,Centro)(9,Ensanche)(16,Periferia)};
      \addplot coordinates {(2,Centro)(3,Ensanche)(4,Periferia)};
    \end{axis}
  \end{tikzpicture}\caption{Horizontal apilado (100 \%)}
  \end{subfigure}

  \begin{subfigure}{0.48\textwidth}
  \begin{tikzpicture}
    \begin{axis}[tfmapiladas, width=\linewidth, height=5.5cm, ylabel={Tramos}, symbolic x coords={Centro,Ensanche,Periferia}, cycle list name=tfmaccesibilidad]
      \addplot coordinates {(Centro,52)(Ensanche,38)(Periferia,25)};
      \addplot coordinates {(Centro,24)(Ensanche,30)(Periferia,28)};
      \addplot coordinates {(Centro,15)(Ensanche,20)(Periferia,27)};
      \addplot coordinates {(Centro,7)(Ensanche,9)(Periferia,16)};
    \end{axis}
  \end{tikzpicture}\caption{Vertical apilado}
  \end{subfigure}\hfill
  \begin{subfigure}{0.48\textwidth}
  \begin{tikzpicture}
    \begin{axis}[tfmlineasv, width=\linewidth, height=5.5cm, xlabel={Ancho libre (m)}, ylabel={Distancia (m)}]
      \addplot coordinates {(2.1,0)(1.9,10)(1.4,20)(1.6,30)(2.0,40)}; \addlegendentry{Acera norte}
      \addplot coordinates {(1.5,0)(1.2,10)(0.9,20)(1.1,30)(1.3,40)}; \addlegendentry{Acera sur}
    \end{axis}
  \end{tikzpicture}\caption{Líneas en vertical (perfil)}
  \end{subfigure}
  \caption{Variantes de barras y líneas.}
\end{figure}

\section{Gráfico circular}
\begin{figure}[htbp]
  \centering
  \graficocircular{42/acc-accesible/Accesible, 27/acc-practicable/Practicable, 19/acc-deficiente/Deficiente, 9/acc-no-accesible/No accesible, 3/acc-sin-datos/Sin datos}
  \hspace{12mm}
  \graficocircular[hueco=10mm, leyenda=true, centro={1\,240\\[-2pt]{\scriptsize\color{tinta-suave}tramos}}]{42/acc-accesible/Accesible, 27/acc-practicable/Practicable, 19/acc-deficiente/Deficiente, 9/acc-no-accesible/No accesible, 3/acc-sin-datos/Sin datos}
  \caption{Distribución de tramos por clase: sectores y anillo (datos de ejemplo).}
\end{figure}

\section{Cronograma}
\begin{figure}[htbp]
  \centering
  \begin{ganttchart}[tfmgantt]{1}{16}
    \gantttitle{Semanas}{16} \\
    \gantttitlelist{1,...,16}{1} \\
    \ganttgroup{Datos}{1}{5} \\
    \ganttbar{Descarga PNOA}{1}{2} \\
    \ganttbar{Anotación}{2}{5} \\
    \ganttmilestone{PEC 2}{5} \\
    \ganttgroup{Modelos}{6}{11} \\
    \ganttbar{Entrenamiento}{6}{9} \\
    \ganttbar{Clasificación}{9}{11} \\
    \ganttmilestone{PEC 3}{11} \\
    \ganttbar{Memoria}{10}{16}
    \ganttlink{elem2}{elem4}
  \end{ganttchart}
  \caption{Planificación del trabajo.}
\end{figure}

\section{Mapa de tramos}
\begin{figure}[htbp]
  \centering
  \begin{tikzpicture}[tfm]
    \fill[rejilla] (0,0) rectangle (2.6,1.6); \fill[rejilla] (3.2,0) rectangle (6,1.6);
    \fill[rejilla] (0,2.2) rectangle (2.6,3.6); \fill[rejilla] (3.2,2.2) rectangle (6,3.6);
    \draw[tramo, acc-accesible] (-0.3,1.9) -- (2.9,1.9);
    \draw[tramo, acc-deficiente] (2.9,1.9) -- (6.3,1.9);
    \draw[tramo, acc-practicable] (2.9,-0.3) -- (2.9,1.9);
    \draw[tramo, acc-no-accesible] (2.9,1.9) -- (2.9,3.9);
  \end{tikzpicture}
  \caption{Tramos coloreados por clase de accesibilidad.}
\end{figure}

\section{Citas}
Cita corta en línea: según los autores, el modelo \enquote{unifica un codificador jerárquico con un decodificador ligero}~\cite{xie2021segformer}\traduccionpropia.

\begin{citalarga}[Xie et al.~\cite{xie2021segformer}, p.~2, traducción propia]
Texto literal de ejemplo de una cita larga, de cuarenta palabras o más, que se compone en un bloque sangrado por ambos lados, con un cuerpo de letra menor, sin comillas y con la referencia al pie, alineada a la derecha y precedida de una raya.
\end{citalarga}

\begin{citanorma}[El subrayado es nuestro.]{Orden TMA/851/2021, art.~5}
Texto literal del artículo citado, copiado sin modificar y con la referencia exacta de la norma, el artículo y el apartado.
\end{citanorma}

\section{Estado del arte}
\begin{table}[htbp]
  \centering
  \caption{Comparativa de trabajos previos (ejemplo de formato). \leyendacumple.}
  \small
  \begin{tabular}{llcccl}
    \toprule
    \cabeceratabla Trabajo & Arquitectura & Aérea & Aceras & Accesib. & Métrica \\
    \midrule
    Long et al.~\cite{long2015fcn} & FCN & \cumple{no} & \cumple{no} & \cumple{no} & mIoU \\
    Ronneberger et al.~\cite{ronneberger2015unet} & U-Net & \cumple{no} & \cumple{no} & \cumple{no} & IoU \\
    Chen et al.~\cite{chen2018deeplab} & DeepLabV3+ & \cumple{no} & \cumple{parcial} & \cumple{no} & mIoU \\
    Xie et al.~\cite{xie2021segformer} & SegFormer & \cumple{no} & \cumple{parcial} & \cumple{na} & mIoU \\
    \filapropia Este trabajo & SegFormer-B2 & \cumple{si} & \cumple{si} & \cumple{si} & mIoU, MAE \\
    \bottomrule
  \end{tabular}
\end{table}

\begin{figure}[htbp]
  \centering
  \lineatemporal[desde=2014, hasta=2024]{2015/FCN/Long et al.~\cite{long2015fcn}, 2017/PSPNet/Zhao et al., 2018/DeepLabV3+/Chen et al.~\cite{chen2018deeplab}, 2021/SegFormer/Xie et al.~\cite{xie2021segformer}, 2023/SAM/Kirillov et al.~\cite{kirillov2023sam}}
  \caption{Evolución de las arquitecturas de segmentación semántica.}
\end{figure}

\begin{figure}[htbp]
  \centering
  \begin{tikzpicture}[tfm, node distance=5mm and 8mm]
    \node[prismacaja] (id) {Registros identificados en Scopus y WoS\\(n = XX)};
    \node[prismacaja, below=of id] (cr) {Registros cribados por título y resumen\\(n = XX)};
    \node[prismacaja, below=of cr] (el) {Textos completos evaluados\\(n = XX)};
    \node[prismacaja, below=of el] (in) {Estudios incluidos\\(n = XX)};
    \node[prismaexcluido, right=of id] (x1) {Duplicados eliminados (n = XX)};
    \node[prismaexcluido, right=of cr] (x2) {Excluidos: no usan imagen aérea (n = XX)};
    \node[prismaexcluido, right=of el] (x3) {Excluidos: sin métricas comparables (n = XX)};
    \foreach \a/\b in {id/cr, cr/el, el/in} \draw[flecha] (\a) -- (\b);
    \foreach \a/\b in {id/x1, cr/x2, el/x3} \draw[flecha] (\a) -- (\b);
    \node[prismafase, minimum width=10mm] at ([xshift=-6mm]id.west) {Identif.};
    \node[prismafase, minimum width=21mm] at ([xshift=-6mm,yshift=-7.5mm]cr.west) {Cribado};
    \node[prismafase, minimum width=10mm] at ([xshift=-6mm]in.west) {Inclusión};
  \end{tikzpicture}
  \caption{Selección de artículos (diagrama PRISMA).}
\end{figure}

\begin{figure}[htbp]
  \centering
  \begin{tikzpicture}
    \begin{axis}[tfmdispersion, width=0.7\textwidth, height=0.42\textwidth, xlabel={Parámetros (M)}, ylabel={mIoU (\%)}, xmin=0, xmax=160, ymin=60, ymax=90]
      \addplot[serie-1, mark=*] table[x=p, y=m, meta=etiqueta] {
        p   m    etiqueta
        134 65   FCN-8s
        31  72   U-Net
        44  79   DeepLabV3+
        27  81   SegFormer-B2
      };
    \end{axis}
  \end{tikzpicture}
  \caption{Precisión frente a tamaño del modelo (datos de ejemplo).}
\end{figure}

\section{Real frente a predicho}
\begin{figure}[htbp]
  \centering
  \begin{tikzpicture}
    \begin{axis}[tfmrealpredicho, width=0.5\textwidth, xmin=0, xmax=4, ymin=0, ymax=4,
        xlabel={Ancho real (m)}, ylabel={Ancho predicho (m)}]
      \fill[bandatolerancia] (axis cs:0,0.3) -- (axis cs:3.7,4) -- (axis cs:4,4) -- (axis cs:4,3.7) -- (axis cs:0.3,0) -- (axis cs:0,0) -- cycle;
      \draw[diagonal] (axis cs:0,0) -- (axis cs:4,4);
      \addplot coordinates {(1.2,1.3)(1.8,1.6)(2.4,2.5)(3.1,2.6)(0.9,1.4)(2.0,2.1)(2.8,2.9)(1.5,1.4)};
    \end{axis}
  \end{tikzpicture}
  \caption{Ancho de acera real frente a predicho; banda de $\pm 0{,}3$~m.}
\end{figure}

\tolerancia{0.3}
\begin{table}[htbp]
  \centering
  \caption{Comparación por tramo (datos de ejemplo).}
  \begin{tabular}{lrrrll}
    \toprule
    \cabeceratabla Tramo & Real (m) & Pred. (m) & Error (m) & Clase real $\to$ pred. & \\
    \midrule
    C/ Mayor 1   & 1,20 & 1,30 & \errorval{0.10}  & 3 $\to$ 3 & \comparaclase{3}{3} \\
    C/ Mayor 2   & 3,10 & 2,60 & \errorval{-0.50} & 1 $\to$ 2 & \comparaclase{1}{2} \\
    Av. Sur 4    & 0,90 & 1,40 & \errorval{0.50}  & 4 $\to$ 3 & \comparaclase{4}{3} \\
    \midrule
    \multicolumn{3}{l}{MAE / RMSE} & 0,37 / 0,41 & \multicolumn{2}{l}{Acierto de clase: 33\,\%} \\
    \bottomrule
  \end{tabular}
\end{table}

\begin{figure}[htbp]
  \centering
  \begin{tikzpicture}[tfm]
    \fill[rejilla] (0,0) rectangle (2.6,1.6); \fill[rejilla] (3.2,0) rectangle (6,1.6);
    \draw[tramoacierto] (-0.3,1.9) -- (2.9,1.9);
    \draw[tramosobre] (2.9,1.9) -- (6.3,1.9);
    \draw[tramosub] (2.9,-0.3) -- (2.9,1.9);
  \end{tikzpicture}
  \caption{Mapa de errores de clase por tramo.}
\end{figure}

\section{Notas y aclaraciones}
Texto de la memoria con una nota al pie\footnote{Las notas al pie siguen siendo la forma preferente para comentarios breves.}.

\postit[Víctor]{Comprobar con el tutor si incluimos los tramos peatonales interiores.}

\aclaracion{En este trabajo, «tramo» es el segmento de acera entre dos cruces consecutivos.}

\begin{table}[htbp]
  \centering
  \caption{Tabla con nota.}
  \begin{tabular}{lr}\toprule \cabeceratabla Clase & Tramos \\ \midrule Accesible & 520 \\ \bottomrule\end{tabular}
  \notatabla{Se excluyen los tramos sin ortofoto de 2025.}
\end{table}

\section{Listas}
Viñetas por niveles:
\begin{itemize}
  \item Imagen aérea de alta resolución (PNOA, 25~cm/píxel).
  \item Anotación de clases:
  \begin{itemize}
    \item calzada y acera;
    \item pasos de peatones;
    \begin{itemize}\item con rebaje;\item sin rebaje.\end{itemize}
    \item obstáculos.
  \end{itemize}
  \item Validación en campo.
\end{itemize}

Numeradas:
\begin{enumerate}
  \item Descargar las ortofotos.
  \item Preparar el conjunto de datos:
  \begin{enumerate}
    \item teselar a $512\times512$;
    \item dividir en entrenamiento y validación.
  \end{enumerate}
  \item Entrenar el modelo.\label{it:entrenar}
\end{enumerate}
El paso~\ref{it:entrenar} es el más costoso.

\begin{description}
  \item[Tramo] Segmento de acera entre dos cruces consecutivos.
  \item[Ancho libre] Distancia mínima sin obstáculos medida perpendicularmente al eje.
\end{description}

Objetivos:
\begin{codigos}[O]
  \item Generar un conjunto de datos anotado.\label{obj:datos}
  \item Entrenar un modelo de segmentación.
  \begin{codigos}\item Comparar arquitecturas.\end{codigos}
  \item Clasificar los tramos por accesibilidad.
\end{codigos}
El objetivo~\ref{obj:datos} condiciona el resto.

\begin{codigos}[H]
  \item La segmentación estima el ancho libre con un error $< 0{,}3$~m.
\end{codigos}

\begin{pasos}
  \item Segmentar la ortofoto.
  \item Vectorizar la máscara.
  \item Calcular métricas por tramo.
\end{pasos}

\begin{ventajas}
  \item Coste por kilómetro muy inferior a la inspección in situ.
  \item Actualizable con cada vuelo del PNOA.
\end{ventajas}
\begin{inconvenientes}
  \item No detecta pendientes ni rebajes finos.
\end{inconvenientes}

\begin{comprobacion}
  \item[\hecho] Dataset anotado.
  \item[\parcial] Modelo entrenado.
  \item[\porhacer] Validación en campo.
\end{comprobacion}

En línea: las clases son \begin{enlinea}\item calzada \item acera \item obstáculo\end{enlinea}.

\section{Código}
\begin{lstlisting}[language=Python, caption={Definición del modelo.}]
import segmentation_models_pytorch as smp
model = smp.Unet(encoder_name="resnet34", classes=8)  # 8 clases
\end{lstlisting}

\printbibliography[heading=bibintoc]
\end{document}
